\chapter{Introduction}

On September 07, 2017, the American credit bureau Equifax announced that they have been a victim of a data breach attack that resulted in up to 148 million personal records being leaked \cite{report}.
\newline
Despite the incident being reported in September, it happened in May 13, 2027, where the attackers exploited a software vulnerability in Equifax's systems that Equifax themselves have not patched, despite the vulnerability being public since May 07, 2017 \cite{arstechnicabug}. 
\newline
Due to this know vulnerability, the attackers were capable of gaining remote access to Equifax's network and found a file containing employee credentials. Using these credentials, the attackers performed up to 9.000 database searches within Equifax's databases to retrieve personal records from customers. On July 29, 2017, Equifax updated it's security certificates and immediately noticed suspicious network traffic and realized that they have been breached. The attack was unnoticed for 76 days \cite{report}.
\newline
As previously mentioned, this breach resulted in up to 148 million personal records being leaked that included full names, addresses, birth dates and social security numbers. Not only did this breach affect their customers, it also exposed Equifax's lack of seriousness in cybersecurity practices as they didn't apply the patch to the software vulnerability when it was available or that they updated their security certificate 19 months after it had already expired. According to the report of the U.S. Committee on Oversight and Government Reform, the attack was preventable \cite{report}.
\newline
The goal of this task is to analyze the incident from a cybersecurity architecture perspective. Using the given evidence of the report, we aim to reconstruct the current-state architecture that was used during the time of the incident and how it escalated from a "mere" vulnerability to a business loss. Additionally, we also aim to propose an architecture that could've been used by Equifax to mitigate the incident. 
